\subsection{Свойства функций, имеющих пределы: арифметические свойства, локальная ограниченность и теорема о сохранении знака.}

\textbf{Теорема (О сохранении знака).} 
Если $\lim_{x\to a} f(x) = A \neq 0$, то существует проколотая окрестность точки $a$, в которой знак функции совпадает со знаком предела:
\[ \exists \delta > 0\; \forall x: 0 < |x - a| < \delta \implies \sign f(x) = \sign A \]

\begin{tcolorbox}[title=Доказательство]
	Возьмем $\varepsilon = \frac{|A|}{2} > 0$.
	По определению предела $\exists \delta > 0$, такое что для всех $0 < |x - a| < \delta$:
	\[ |f(x) - A| < \varepsilon = \frac{|A|}{2} \]
	Раскроем модуль:
	\[ A - \frac{|A|}{2} < f(x) < A + \frac{|A|}{2} \]
	Рассмотрим два случая:
	\begin{enumerate}
		\item Пусть $A > 0$. Тогда $|A| = A$.
		\[ f(x) > A - \frac{A}{2} = \frac{A}{2} > 0 \implies \sign f(x) = 1 = \sign A \]
		\item Пусть $A < 0$. Тогда $|A| = -A$.
		\[ f(x) < A + \frac{-A}{2} = \frac{A}{2} < 0 \implies \sign f(x) = -1 = \sign A \]
	\end{enumerate}
\end{tcolorbox}

\textbf{Теорема (О локальной ограниченности).}
Если существует конечный предел $\lim_{x \to a} f(x) = A$, то функция $f(x)$ ограничена в некоторой проколотой окрестности точки $a$.

\begin{tcolorbox}[title=Доказательство]
	Пусть $\varepsilon = 1$.
	По определению предела $\exists \delta > 0$, такое что $\forall x: 0 < |x - a| < \delta$:
	\[ |f(x) - A| < 1 \]
	Используем неравенство треугольника $|a| - |b| \le |a+b|$:
	\[ |f(x)| = |(f(x) - A) + A| \leq |f(x) - A| + |A| < 1 + |A| \]
	Обозначим $M = |A| + 1$. Тогда $|f(x)| < M$.
	Функция ограничена.
\end{tcolorbox}

\textbf{Теорема (Арифметические свойства).}
Пусть $f(x) \xrightarrow{x\to a} A$ и $g(x) \xrightarrow{x\to a} B$. Тогда:
\begin{enumerate}
	\item $\lim_{x\to a} (f(x) \pm g(x)) = A \pm B$
	\item $\lim_{x\to a} (f(x) \cdot g(x)) = A \cdot B$
	\item Если $B \neq 0$, то $\lim_{x\to a} \frac{f(x)}{g(x)} = \frac{A}{B}$
\end{enumerate}

\begin{tcolorbox}[title=Доказательство (через Гейне)]
	Используем эквивалентность определений по Коши и по Гейне.
	
	Пусть $\{x_n\}$ — произвольная последовательность такая, что $x_n \in \mathring{U}_a$ и $x_n \to a$.
	Так как пределы функций существуют, то по определению Гейне соответствующие числовые последовательности значений сходятся:
	\[ f(x_n) \to A \quad \text{и} \quad g(x_n) \to B \quad \text{при } n \to \infty \]
	
	Тогда сводим задачу к свойствам пределов числовых последовательностей:
	\begin{enumerate}
		\item $f(x_n) \pm g(x_n) \to A \pm B$
		\item $f(x_n) \cdot g(x_n) \to A \cdot B$
		\item Для частного:
		Так как $g(x_n) \to B$ и $B \neq 0$, то начиная с некоторого номера $N$ выполняется $g(x_n) \neq 0$. Следовательно, частное определено и:
		\[ \frac{f(x_n)}{g(x_n)} \to \frac{A}{B} \]
	\end{enumerate}
	
	Так как это верно для любой последовательности $\{x_n\}$, то по определению Гейне пределы функций существуют и равны результатам арифметических операций над $A$ и $B$.
\end{tcolorbox}
